\chapter*{Abstract}

Hand exoskeletons for rehabilitation must satisfy demanding requirements in terms of mechanical compliance, dynamic accuracy, and operational safety, while interacting directly with the human body. Developing and validating such systems through iterative hardware prototyping alone is time-consuming, costly, and potentially risky. This thesis addresses these challenges by developing a comprehensive Digital Twin of a hand exoskeleton that integrates physics-based modeling, closed-loop control, functional safety, and model-based fault detection within a unified ROS 2 framework.
 
The exoskeleton mechanism, characterized by multiple closed kinematic chains driven by a single actuated joint, was modeled parametrically in Fusion 360 and translated into a URDF description compatible with standard robotic simulation tools. Because the URDF formalism does not support closed-loop topologies, the kinematic chains were deliberately opened and the loop-closure conditions were reimposed at the software level through a constraint-based formulation built on the Pinocchio rigid-body dynamics library. This approach yields the implicit mapping $q = q(\theta)$ that associates each crank angle with a unique, kinematically admissible configuration of the full mechanism.
 
A reduced-order dynamic model was derived by projecting the constrained multibody dynamics onto the single actuated degree of freedom. The projection condenses the coupled inertial, Coriolis, gravitational, friction, and passive biomechanical effects into a scalar nonlinear equation that is evaluated numerically at each simulation step, preserving the nonlinear physics of the original system while ensuring deterministic execution at 200 Hz. An extended version of the model incorporates the finger phalanges and a Fung-type exponential stiffness model for the passive tissue response.
 
A cascade control architecture was implemented on top of the dynamic model. The outer admittance loop generates a compliant motion reference by simulating a virtual mass-damper-spring system driven by the user's interaction force, while the inner trajectory loop tracks this reference through a PD controller with comprehensive model-based feed-forward compensation. Closed-loop validation through step and sinusoidal response experiments confirmed stable, overshoot-free transients with sub-degree tracking accuracy.
 
A three-layer functional safety architecture was developed to protect the user against software and sensor faults. The enforcement layer — an ExoBridge gateway node — mediates all control signals, applying mode-dependent torque and velocity limits and maintaining per-channel communication watchdogs. The detection layer — a FaultMonitorMode plugin — evaluates the system signals against configurable thresholds through a debounced, multi-level escalation logic. The governance layer — a plugin-based C++ state machine — manages four safety states (Fault Monitor, Compliant Mode, Torque Limit Mode, Safe Stop) with a well-defined transition graph, fault latching, bridge liveness and mode coherence supervision, and an optional automatic downgrade mechanism. A dedicated fault injection framework, supporting six fault types across four signal channels, enables systematic and reproducible validation of the safety mechanisms.
 
A model-based fault detection module, comprising a Luenberger state observer, an admittance inversion estimator, a generalized momentum observer, and a rate-of-change guard, was implemented to explore observer-based sensor health monitoring. The four layers provide complementary coverage of the fault space with detection latencies ranging from 5 ms to 1.5 s. The module is fully operational within the Digital Twin but has not yet been integrated into the safety decision chain; its connection to the state machine through a reserved Sensor Degraded Mode state is identified as future work.
 
The hardware prototype of the exoskeleton, including the HT1225 brushless DC actuators, HT-PTZ field-oriented drive boards, load cell amplifiers, and finger encoders, was comprehensively documented, establishing the interface specification for future hardware-in-the-loop integration.
 
The thesis demonstrates that a Digital Twin approach, combining constrained multibody dynamics, compliant control, layered functional safety, and model-based fault detection, can effectively support the complete development cycle of a hand exoskeleton for rehabilitation, from early-stage design to pre-deployment safety verification.